\section{Transient problems: heat pulses and thermal shock}\label{sec:transient}
Steady-state dispersion analysis tells us how a monochromatic wave propagates. Transient problems test the \emph{fronts}: finite speed, jumps and their decay. Two benchmark problems are used here because their Laplace-domain solutions are closed-form and they can be inverted numerically to any accuracy; they also serve as test cases for numerical codes (Section~\ref{sec:methods}).

\subsection{Uncoupled heat pulse in a half-space}
Consider the semi-infinite rod $x>0$, initially at $\theta=0$, with prescribed boundary flux $q(0,t)=q_0H(t)$ and no coupling. For $G(s)=k(1+\tau_Ts)/(1+\tau_qs)$ with $\tau_T=0$ (Lord--Shulman), $\tau_q=\tau$, the solution is
\begin{equation}
\hat\theta(x,s)=\frac{q_0\,(1+\tau s)}{k\,s\,m(s)}\,e^{-x\,m(s)},\qquad m(s)=\sqrt{\frac{s(1+\tau s)}{\alpha}},\label{eq:pulse}
\end{equation}
and $\tau=0$ gives Fourier's solution $\theta=(q_0/k)\bigl[2\sqrt{\alpha t/\pi}\,e^{-x^2/4\alpha t}-x\,\mathrm{erfc}(x/2\sqrt{\alpha t})\bigr]$. Large-$s$ asymptotics, $m(s)\simeq s/c_T+1/(2\ell)$ with $c_T=\sqrt{\alpha/\tau}$ and $\ell=\sqrt{\alpha\tau}=c_T\tau$, give the front structure: the signal is exactly zero ahead of $x=c_Tt$; at the front it jumps by
\begin{equation}
\bigl[\theta\bigr]_{x=c_Tt}=\frac{q_0}{\rho c\,c_T}\,e^{-x/2\ell},
\end{equation}
decaying exponentially with distance, and at the boundary the temperature jumps from zero to $q_0/\rho c c_T$ at $t=0^+$, instead of growing as $\sqrt t$. The Cattaneo curves of \Cref{fig:pulse} were obtained by numerical Talbot inversion \citep{Talbot1979} of \eqref{eq:pulse} (with the exact zero set ahead of the front); the ripples near the front are the Gibbs-type error of the inversion at a discontinuity. At late times the two solutions coincide (the ratio at $t=40\tau$ differs from one by less than one percent), which is the sense in which the Cattaneo model is a high-frequency correction to Fourier's law.

\begin{figure}[t]\centering
\includegraphics[width=\textwidth]{figs/fig_heatpulse.pdf}
\caption{Uncoupled heat pulse from a prescribed boundary flux. Left: temperature history at $x=1.5\ell$. Right: profile at $t=2\tau$. The Lord--Shulman curve is exactly zero beyond the front (dotted); the classical solution has no front. At early times the Lord--Shulman surface temperature exceeds the classical one for this boundary condition.}\label{fig:pulse}
\end{figure}

\paragraph{Boundary conditions decide the sign.} Statements such as ``second sound lowers the surface temperature'' or ``raises'' it are meaningless without the boundary condition. With a prescribed flux, the hyperbolic model gives a larger surface temperature than Fourier's law at early times in this normalisation (\Cref{fig:pulse}); with a prescribed surface temperature, it is the flux that jumps, with an overshoot relative to Fourier's. Laser-heating models therefore require a careful statement of how the heat input is represented \citep{Li2019}.

\subsection{Thermal shock of a coupled half-space}
Next consider the half-space $x>0$ with the boundary at temperature $\theta_0H(t)$ and traction-free, coupled thermoelastic one-dimensional problem. In the Laplace domain, with fields $\propto e^{-m x}$ the characteristic equation is \eqref{eq:disp} with $K=-m^2$, giving two decay rates $m_1(s)$, $m_2(s)$ ($\Real m>0$). In units $\rho=c=c_1=k=T_0=1$ (so that lengths are in $c_1/\omega_*$ and time in $1/\omega_*$), with $\beta=\sqrt\varepsilon$, the displacement and stress amplitudes associated with a temperature amplitude $A_j$ are
\begin{equation}
U_j=\frac{\beta\,m_j}{s^2-m_j^2}A_j,\qquad \sigma_j=-m_jU_j-\beta A_j
\end{equation}
and the boundary conditions $\sum_jA_j=\theta_0/s$, $\sum_j\sigma_j(0)=0$ determine $A_1,A_2$; the second condition reduces to $\sum_jA_j\,s^2/(s^2-m_j^2)=0$. We verified that the resulting transform satisfies $s\hat\theta(0,s)=1$ and $\hat\sigma(0,s)=0$ to machine precision, and invert numerically. \Cref{fig:shock} shows temperature and stress at a point $x=2c_1/\omega_*$ for the classical, Lord--Shulman and DPL laws ($\varepsilon=0.1$, $\tau_q=0.25/\omega_*$, DPL $\tau_T=\tau_q/2$).

\begin{figure}[t]\centering
\includegraphics[width=\textwidth]{figs/fig_shock.pdf}
\caption{Thermal shock of a coupled half-space (boundary temperature step, traction-free surface), recorded at $x=2c_1/\omega_*$. Dotted line: arrival of the elastic wave front ($x=c_1t$). The Lord--Shulman thermal front reaches the point at $t=x/c_T=1$; the stress develops a negative (compressive) pulse and jumps at the elastic front. Small oscillations are inversion error at discontinuities.}\label{fig:shock}
\end{figure}

The following qualitative features are robust. (i) The thermal part of the response is almost identical in all three models after the first elastic transit, because the differences between heat laws are confined to the first moments after loading; the long-time evolution is governed by the classical diffusion. (ii) The stress is largest, and most sensitive to the heat law, at the elastic front, where the Lord--Shulman model gives about 20\,\% larger compressive peak than the classical one in this example (peak values $-0.30$ against $-0.25$, in units of $\rho c_1^2$). (iii) Before the elastic front arrives the classical solution already carries stress, because the thermal body force acts instantly everywhere; in Lord--Shulman the stress is essentially zero until the thermal front arrives at $t=1$ and then rises more steeply. The numerical amplitudes depend on the dimensionless parameters and should not be extrapolated to real materials without recomputation.

\subsection{Lessons for laser heating and for numerical codes}
Short-pulse laser heating in metals is the main application driving hyperbolic models; analyses of bilayers and films under pulsed heating combine fractional and memory-dependent laws \citep{Li2019}. For numerical codes, the benchmark problems above provide a stringent check, because a code that smears the front also misses the stress peak. The classical statement of the thermal shock problem is due to Danilovskaya \citep{Danilovskaya1950}, and thermal-shock solutions of the coupled problem are a standard test in texts \citep{Boley1960,Nowacki1975,HetnarskiIgnaczak1999}.

\begin{summarybox}[Summary]
The heat-pulse problem shows a finite front for hyperbolic laws (LS, MGT, GN-II) and a Gaussian tail for Fourier. The front carries a stress peak whose height is heat-law dependent (about 20\% between Fourier and LS for the parameters used). Thermal shock combines the elastic front with diffusive wake; a numerical code that smears the front also misses the stress peak---the benchmark of Section~\ref{sec:transient} tests exactly this.
\end{summarybox}
